\documentclass[aspectratio=169, 14pt]{beamer}
\usepackage{beamer_preamble}

\begin{document}

% ==== Title ===============================================================
\titleslide{Lesson 4-5 \textperiodcentered\ Period 1 (Quick)}{Solving Rational Equations + Chemistry DOK-3 (Q\#2/Q\#6/Q\#7 Prep)}

% ==== Do Now ==============================================================
\begin{frame}{Do Now --- Critique \& Explain}{Algebra 2 \textperiodcentered\ Lesson 4-5 \textperiodcentered\ Single-Period Quick}
\phasebadges{1}{5}{bridge from 4-4}
\vspace{0.2cm}
\begin{projcard}{Do Now \textperiodcentered\ Nicky vs.\ Tavon}{slidegold}
\small
Both students solve \(\dfrac{1}{2}x + \dfrac{2}{5} = \dfrac{9}{10}\) --- Nicky subtracts fractions; Tavon multiplies both sides by 10.

\vspace{0.2cm}
\begin{enumerate}
  \setlength{\itemsep}{2pt}
  \item[(A)] What are the advantages/disadvantages of each strategy?
  \item[(B)] Are both methods correct?
  \item[(C)] Why did Tavon choose to multiply by 10?
\end{enumerate}
\end{projcard}
\end{frame}

% ==== Launch ==============================================================
\begin{frame}{Launch --- Rules + Example 1 (Solve)}{Algebra 2 \textperiodcentered\ Lesson 4-5 \textperiodcentered\ Single-Period Quick}
\phasebadges{1--2}{13}{teacher models Ex 1 + Ex 3}
\vspace{0.1cm}
\begin{projcard}{Rules --- Solving Rational Equations}{slidegreen}
\small
\begin{enumerate}
  \setlength{\itemsep}{1pt}
  \item Find the LCD of all fractions.
  \item Multiply BOTH SIDES by the LCD to clear fractions.
  \item Solve the resulting polynomial equation.
  \item CHECK each candidate in the ORIGINAL equation --- any candidate that makes a denominator zero is \textbf{extraneous}.
\end{enumerate}
\vspace{0.1cm}
Domain restrictions are set BEFORE clearing denominators.
\end{projcard}
\end{frame}

\begin{frame}{Launch --- Example 3 (Extraneous Solution)}{Algebra 2 \textperiodcentered\ Lesson 4-5 \textperiodcentered\ Single-Period Quick}
\phasebadges{1--2}{13}{teacher models --- COMPRESSED 2 examples}
\vspace{0.15cm}
\begin{projcard}{Key idea: why candidates can be extraneous}{slidegreen}
\small
Multiplying both sides by an expression containing \(x\) can \textbf{introduce} values that are not in the domain of the original equation.

\vspace{0.1cm}
Example 3 pattern:
\begin{enumerate}
  \setlength{\itemsep}{1pt}
  \item State domain restrictions \textbf{first} (values that make any denominator zero).
  \item Solve the cleared equation --- collect candidate solutions.
  \item Test each candidate against the restrictions. Extraneous = throw it out.
\end{enumerate}

\vspace{0.1cm}
\textcolor{slideaccent}{\textbf{Extraneous-solution trap = \#1 error on this lesson's assessment items!}}
\end{projcard}
\end{frame}

% ==== Explore =============================================================
\begin{frame}{Explore --- Think-Pair-Share (32 min)}{Algebra 2 \textperiodcentered\ Lesson 4-5 \textperiodcentered\ Single-Period Quick}
\phasebadges{2--3}{32}{student work}
\small\textcolor{slidegray}{6 items in order. Plan \(\sim\)10 min for Practice \#33 (Chemistry). Wed F: skip \#14 + \#25.}

\vspace{0.1cm}
\begin{projcard}{Try It 1 \textperiodcentered\ Solve a rational equation}{slideblue}
What is the solution to each equation? \quad
\(\dfrac{2}{x+5} = 4\) \quad \(\dfrac{1}{x-7} = 2\)
\end{projcard}
\vspace{-0.12cm}
\begin{projcard}{Try It 3 \textperiodcentered\ Identify extraneous solution}{slideblue}
\(\dfrac{1}{x+2} + \dfrac{1}{x-2} = \dfrac{4}{(x+2)(x-2)}\) \quad What is the solution?
\end{projcard}
\vspace{-0.12cm}
\begin{projcard}{Practice \#10 \textperiodcentered\ Error analysis (Miranda)}{slideblue}
Describe and correct Miranda's error. What is the correct solution?
\end{projcard}
\end{frame}

\begin{frame}{Explore (cont.) --- Work-Rate + DOK-3}{Algebra 2 \textperiodcentered\ Lesson 4-5 \textperiodcentered\ Single-Period Quick}
\phasebadges{2--3}{32}{student work}
\small\textcolor{slidegray}{Continue in order. Practice \#33 is the big one --- plan \(\sim\)10 min.}

\vspace{0.15cm}
\begin{projcard}{Practice \#14 \textperiodcentered\ Extraneous trap}{slideblue}
Solve \(\dfrac{x^2-7x-18}{x+2} = x-9\). Explain what is unique about the solution.
\end{projcard}
\vspace{-0.12cm}
\begin{projcard}{Practice \#25 \textperiodcentered\ Work-rate (Kenji + Oscar)}{slideblue}
Setup: \(\tfrac{1}{a} + \tfrac{1}{b} = \tfrac{1}{t}\). How long together?
\end{projcard}
\vspace{-0.12cm}
\begin{projcard}{Practice \#33 \textperiodcentered\ PERFORMANCE TASK (\(\bigstar\))}{slideaccent}
Chemistry mixture: set \(f(x) = \dfrac{50(0.02)+0.06x}{50+x}\) equal to target concentration. Solve for \(x\). Verify with TABLE.
\end{projcard}
\end{frame}

% ==== Share / Summary =====================================================
\begin{frame}{Share / Summary}{Algebra 2 \textperiodcentered\ Lesson 4-5 \textperiodcentered\ Single-Period Quick}
\phasebadges{2}{5}{DOK-3 reveal + LEHS bridge}
\vspace{0.2cm}
\begin{projcard}{Sentence frame \textperiodcentered\ Chemistry (\#33)}{slideblue}
\small
``I set up \(f(x) = \dfrac{50\cdot0.02 + 0.06x}{50+x} = 0.05\). After clearing the denominator I got \underline{\hspace{1.2cm}}. Solving for \(x\) gave \(x =\) \underline{\hspace{0.8cm}} gallons. I verified with a calculator TABLE by entering \(y_1 =\) \underline{\hspace{1cm}} and \(y_2 =\) \underline{\hspace{1cm}}.''
\end{projcard}

\vspace{0.15cm}
\begin{projcard}{Bridge to Topic 4 LEHS Assessment}{slideaccent}
\small
\textbf{Q\#2}: domain-restriction reasoning \textperiodcentered\
\textbf{Q\#6}: solve rational equation (same moves as \#33) \textperiodcentered\
\textbf{Q\#7}: work-rate reasoning.
You have now rehearsed all three.
\end{projcard}
\end{frame}

% ==== Exit Ticket =========================================================
\begin{frame}{Exit Ticket --- Summary Recap}{Algebra 2 \textperiodcentered\ Lesson 4-5 \textperiodcentered\ Single-Period Quick}
\phasebadges{1--2}{3}{not graded}
\vspace{0.2cm}
\begin{projcard}{Complete on your exit ticket}{slidegold}
\begin{enumerate}
  \setlength{\itemsep}{0.25cm}
  \item After clearing denominators I must check each candidate in the \underline{\hspace{1.2cm}} equation to catch \underline{\hspace{1.2cm}} solutions.
  \item In Practice \#33 (chemistry), the equation came from setting \underline{\hspace{2.2cm}} equal to the target concentration 0.05.
  \item One biggest thing I learned today: \underline{\hspace{4cm}}.
\end{enumerate}
\end{projcard}
\end{frame}

\end{document}
